% ========================================================
%  习近平新时代中国特色社会主义思想概论  期末全真模拟试卷（含参考答案）
%  覆盖范围：导论、第一至八章、第十至十二章 核心考点
% ========================================================
\documentclass[12pt, a4paper]{article}

% ---- 中文支持 ----
\usepackage[UTF8]{ctex}

% ---- 数学 ----
\usepackage{amsmath, amssymb, amsthm}
\usepackage{mathtools}

% ---- 页面布局 ----
\usepackage[top=2.5cm, bottom=2.5cm, left=2.8cm, right=2.8cm]{geometry}

% ---- 表格 ----
\usepackage{booktabs, array, multirow, makecell}
\usepackage{tabularx}

% ---- 页眉页脚 ----
\usepackage{fancyhdr}
\usepackage{lastpage}
\setlength{\headheight}{14pt}
\pagestyle{fancy}
\fancyhf{}
\renewcommand{\headrulewidth}{0.6pt}
\lhead{\small 习近平新时代中国特色社会主义思想概论·Claude·A卷}
\rhead{\small 共 \pageref{LastPage} 页}
\cfoot{\small 第 \thepage\ 页}

% ---- 计数器与样式 ----
\usepackage{enumitem}
\usepackage{xcolor}
\usepackage{tcolorbox}
\tcbuselibrary{skins, breakable}

% ---- 填空线 ----
\newcommand{\blank}[1]{\underline{\hspace{#1}}}

% ---- 选择题选项列表 ----
\newlist{choices}{enumerate}{1}
\setlist[choices]{label=\textbf{\Alph*.}, itemsep=1pt, topsep=2pt, parsep=0pt, leftmargin=2.4em}

% ---- 答案盒子 ----
\newtcolorbox{answerbox}[1][]{
  breakable,
  colback=gray!6,
  colframe=gray!50,
  fonttitle=\bfseries\small,
  title=参考答案,
  #1
}

% ---- 题号格式 ----
\newcounter{bigq}
\newcommand{\BigQ}[1]{%
  \stepcounter{bigq}%
  \vspace{6pt}%
  \noindent{\large\bfseries 第\chinese{bigq}大题\quad #1}%
  \vspace{4pt}\par
}

\newcounter{subq}[bigq]
\newcommand{\SubQ}{%
  \stepcounter{subq}%
  \noindent\textbf{（\arabic{subq}）}\quad
}

% 分隔线
\newcommand{\examrule}{\noindent\rule{\linewidth}{0.6pt}}

% ---- 文档开始 ----
\begin{document}

% ============================================================
%  封面
% ============================================================
\begin{center}
  {\LARGE\bfseries 习近平新时代中国特色社会主义思想概论}\\[8pt]
  {\large\bfseries 期末全真模拟测试卷}\\[16pt]
  \begin{tabular}{lll}
    考试时间：90 分钟 & \quad & 满分：100 分 \\[4pt]
    \multicolumn{3}{l}{注意事项：请将答案写在答题纸指定区域，写在本卷上无效。}
  \end{tabular}
  \vspace{4pt}

  \examrule

  \vspace{6pt}
  \begin{tabular}{|c|c|c|c|c|}
    \hline
    \textbf{题号} & 一（单选） & 二（多选） & 三（简答） & \textbf{总分} \\
    \hline
    \textbf{满分} & 30 & 20 & 50 & 100 \\
    \hline
    \textbf{得分} & & & & \\
    \hline
  \end{tabular}
\end{center}

\vspace{10pt}
\examrule
\vspace{6pt}

% ============================================================
%  第一大题：单项选择题
% ============================================================
\BigQ{单项选择题（每题只有一个正确答案，每题 1 分，共 30 分）}

\SubQ 习近平新时代中国特色社会主义思想创立的时代背景，不包括以下哪一方面？
\begin{choices}
\item 世界百年未有之大变局加速演进
\item 中华民族伟大复兴进入关键时期
\item 中国式现代化全面推进拓展
\item 全面建成小康社会的历史任务已经完成
\end{choices}

\SubQ "十个明确"在习近平新时代中国特色社会主义思想科学体系中的地位是？
\begin{choices}
\item 是这一思想的世界观和方法论
\item 是新时代坚持和发展中国特色社会主义的基本方略
\item 是这一思想最为核心关键的组成部分，是起支撑作用的"四梁八柱"
\item 全面展示了这一思想的实践成果
\end{choices}

\SubQ 习近平新时代中国特色社会主义思想是"两个结合"的重大成果，其中被称为"魂脉"的是？
\begin{choices}
\item 中华优秀传统文化
\item 马克思主义基本原理
\item 中国具体实际
\item 科学社会主义基本原则
\end{choices}

\SubQ 关于习近平新时代中国特色社会主义思想的历史地位，下列表述不准确的是？
\begin{choices}
\item 是当代中国马克思主义、二十一世纪马克思主义
\item 是中华文化和中国精神的时代精华
\item 实现了马克思主义中国化时代化新的飞跃
\item 实现了马克思主义中国化时代化"第三次历史性飞跃"
\end{choices}

\SubQ "两个确立"具体指的是？
\begin{choices}
\item 确立习近平同志党中央的核心、全党的核心地位；确立习近平新时代中国特色社会主义思想的指导地位
\item 维护党中央权威和集中统一领导；维护党中央和各级党组织的领导地位
\item 坚持和加强党的全面领导；坚持中国特色社会主义道路
\item 中国共产党的领导地位确立；社会主义市场经济体制确立
\end{choices}

\SubQ 科学社会主义基本原则首次在哪部著作中得到全面阐述？
\begin{choices}
\item 《资本论》
\item 《哥达纲领批判》
\item 《共产党宣言》
\item 《反杜林论》
\end{choices}

\SubQ 中国特色社会主义进入新时代，体现的是？
\begin{choices}
\item 我国发展新的历史方位
\item 我国经济发展进入新常态
\item 我国社会主要矛盾已经彻底解决
\item 我国已全面建成社会主义现代化强国
\end{choices}

\SubQ 关于"一个变、两个没有变"，下列表述正确的是？
\begin{choices}
\item 变化的是我国仍处于社会主义初级阶段的基本国情
\item 变化的是我国社会主要矛盾
\item 没有变的是我国社会主要矛盾
\item 没有变的是我国世界最大发展中国家的国际地位已经改变
\end{choices}

\SubQ "道路问题是关系党的事业兴衰成败的第一位的问题"，这一论断强调的核心逻辑是？
\begin{choices}
\item 一个国家的发展速度比道路选择更重要
\item 一个国家实行什么主义，关键要看这个主义能否解决这个国家面临的历史性课题
\item 道路问题只是经济发展中的次要问题
\item 不同国家的发展道路必然趋同
\end{choices}

\SubQ 中国特色社会主义创造的"两个奇迹"是指？
\begin{choices}
\item 经济快速发展奇迹和社会长期稳定奇迹
\item 脱贫攻坚奇迹和乡村振兴奇迹
\item 科技创新奇迹和教育发展奇迹
\item 文化繁荣奇迹和生态文明奇迹
\end{choices}

\SubQ 下列哪一项不属于中国式现代化的中国特色？
\begin{choices}
\item 人口规模巨大的现代化
\item 全体人民共同富裕的现代化
\item 完全照搬西方模式的现代化
\item 人与自然和谐共生的现代化
\end{choices}

\SubQ 推进中国式现代化必须牢牢把握的重大原则不包括？
\begin{choices}
\item 坚持和加强党的全面领导
\item 坚持以人民为中心的发展思想
\item 坚持发扬斗争精神
\item 坚持优先发展重化工业
\end{choices}

\SubQ "团结奋斗是中国共产党、中国人民最显著的精神标识"，这一论述体现的是哪条"必由之路"？
\begin{choices}
\item 中国特色社会主义是实现中华民族伟大复兴的必由之路
\item 团结奋斗是中国人民创造历史伟业的必由之路
\item 自我革命是党永葆生机活力的必由之路
\item 中国式现代化是强国建设的必由之路
\end{choices}

\SubQ 我国当前最大的国情是？
\begin{choices}
\item 社会主义初级阶段
\item 中国共产党的领导
\item 人口规模巨大
\item 区域发展不平衡
\end{choices}

\SubQ 在党的自身优势"四力"中，居于第一位的能力是？
\begin{choices}
\item 思想引领力
\item 群众组织力
\item 社会号召力
\item 政治领导力
\end{choices}

\SubQ 党的领导的最高原则是？
\begin{choices}
\item 民主集中制
\item 党中央集中统一领导
\item 党内监督
\item 群众路线
\end{choices}

\SubQ 中国共产党的根本政治立场是？
\begin{choices}
\item 党的领导
\item 人民立场
\item 依法治国
\item 改革开放
\end{choices}

\SubQ 关于推进共同富裕，下列表述正确的是？
\begin{choices}
\item 应整齐划一、同步实现
\item 要坚持有底线、渐进式推进
\item 只是经济问题，不涉及政治问题
\item 应优先发展东部，暂缓中西部发展
\end{choices}

\SubQ 全面深化改革的总目标是？
\begin{choices}
\item 实现共同富裕、消除贫富差距
\item 完善和发展中国特色社会主义制度，推进国家治理体系和治理能力现代化
\item 建立社会主义市场经济体制
\item 实现高水平科技自立自强
\end{choices}

\SubQ 统筹推进各领域各方面改革开放，应以什么为突破口？
\begin{choices}
\item 制度建设
\item 经济体制改革
\item 重要领域和关键环节改革
\item 党的建设制度改革
\end{choices}

\SubQ 新发展理念中，主要解决发展不平衡问题的是？
\begin{choices}
\item 创新发展
\item 协调发展
\item 绿色发展
\item 共享发展
\end{choices}

\SubQ "两个毫不动摇"指的是？
\begin{choices}
\item 毫不动摇巩固和发展公有制经济，毫不动摇鼓励、支持、引导非公有制经济发展
\item 毫不动摇扩大开放，毫不动摇推进改革
\item 毫不动摇发展教育，毫不动摇发展科技
\item 毫不动摇维护稳定，毫不动摇促进发展
\end{choices}

\SubQ "科技是第一生产力、人才是第一资源、创新是第一动力"，这一表述强调的是？
\begin{choices}
\item 科技、人才、创新三者相互独立、互不影响
\item 教育优先发展、科技自立自强、人才引领驱动一体推进
\item 经济建设优先于其他各项建设
\item 人才工作可以脱离党的领导
\end{choices}

\SubQ 我们党强调的全过程人民民主，是过程民主和什么的统一？
\begin{choices}
\item 结果民主
\item 成果民主
\item 形式民主
\item 程序民主
\end{choices}

\SubQ 我国的根本政治制度是？
\begin{choices}
\item 中国共产党领导的多党合作和政治协商制度
\item 人民代表大会制度
\item 民族区域自治制度
\item 基层群众自治制度
\end{choices}

\SubQ 在"四个自信"中，被认为是更基础、更广泛、更深沉的力量的是？
\begin{choices}
\item 道路自信
\item 理论自信
\item 制度自信
\item 文化自信
\end{choices}

\SubQ "举旗帜、聚民心、育新人、兴文化、展形象"中，"育新人"主要强调的是？
\begin{choices}
\item 高举马克思主义、中国特色社会主义旗帜
\item 牢牢把握正确舆论导向
\item 坚持立德树人、以文化人，培养时代新人
\item 推进国际传播能力建设
\end{choices}

\SubQ 收入分配制度改革中，发挥基础性作用的分配环节是？
\begin{choices}
\item 初次分配
\item 再分配
\item 第三次分配
\item 转移支付
\end{choices}

\SubQ 党的十八大将生态文明建设纳入了什么总体布局？
\begin{choices}
\item "四个全面"战略布局
\item "五位一体"总体布局
\item "三大攻坚战"布局
\item "六稳六保"布局
\end{choices}

\SubQ "绿水青山就是金山银山"科学内涵揭示的道理是？
\begin{choices}
\item 经济发展必须以牺牲生态环境为代价
\item 保护生态环境就是保护生产力，改善生态环境就是发展生产力
\item 生态环境保护与经济发展毫无关联
\item 生态文明建设只是地方政府的责任
\end{choices}

\vspace{10pt}
\examrule

% ============================================================
%  第二大题：多项选择题
% ============================================================
\BigQ{多项选择题（每题有两个或两个以上正确答案，多选、少选、错选均不得分；每题 2 分，共 20 分）}

\SubQ 习近平新时代中国特色社会主义思想创立的时代背景包括哪些方面？
\begin{choices}
\item 世界百年未有之大变局加速演进
\item 中华民族伟大复兴进入关键时期
\item 中国式现代化全面推进拓展
\item 科学社会主义在21世纪的中国焕发新的蓬勃生机
\item 中国共产党在革命性锻造中更加坚强有力
\end{choices}

\SubQ 关于"一个变、两个没有变"，下列说法正确的是？
\begin{choices}
\item 我国社会主要矛盾已经发生变化
\item 我国仍处于并将长期处于社会主义初级阶段的基本国情没有变
\item 我国是世界最大发展中国家的国际地位没有变
\item 我国社会主要矛盾的变化，改变了我们对社会主义所处历史阶段的判断
\item 我国已经全面实现社会主义现代化
\end{choices}

\SubQ 中国式现代化的中国特色主要体现在哪些方面？
\begin{choices}
\item 人口规模巨大的现代化
\item 全体人民共同富裕的现代化
\item 物质文明和精神文明相协调的现代化
\item 人与自然和谐共生的现代化
\item 走和平发展道路的现代化
\end{choices}

\SubQ 推进中国式现代化必须牢牢把握的重大原则包括？
\begin{choices}
\item 坚持和加强党的全面领导
\item 坚持中国特色社会主义道路
\item 坚持以人民为中心的发展思想
\item 坚持深化改革开放
\item 坚持发扬斗争精神
\end{choices}

\SubQ 党的自身优势主要包括哪些方面的能力？
\begin{choices}
\item 政治领导力
\item 思想引领力
\item 群众组织力
\item 社会号召力
\item 经济调控力
\end{choices}

\SubQ 新发展理念的科学内涵主要包括？
\begin{choices}
\item 创新发展
\item 协调发展
\item 绿色发展
\item 开放发展
\item 共享发展
\end{choices}

\SubQ 全过程人民民主体现了哪几个方面的统一？
\begin{choices}
\item 过程民主和成果民主的统一
\item 程序民主和实质民主的统一
\item 直接民主和间接民主的统一
\item 人民民主和国家意志的统一
\item 结果民主和形式民主的统一
\end{choices}

\SubQ 我国的基本政治制度包括？
\begin{choices}
\item 中国共产党领导的多党合作和政治协商制度
\item 民族区域自治制度
\item 基层群众自治制度
\item 人民代表大会制度
\item 民主集中制
\end{choices}

\SubQ 中国特色社会主义文化主要由哪些部分构成？
\begin{choices}
\item 中华优秀传统文化
\item 革命文化
\item 社会主义先进文化
\item 西方现代文明成果
\item 宗教文化
\end{choices}

\SubQ 推进中国式现代化需要正确处理的重大关系包括？
\begin{choices}
\item 顶层设计与实践探索的关系
\item 战略与策略的关系
\item 守正与创新的关系
\item 效率与公平的关系
\item 活力与秩序的关系
\end{choices}

\vspace{10pt}
\examrule

% ============================================================
%  第三大题：简答题
% ============================================================
\BigQ{简答题（请简要论述，只写出结论不进行论述不能得满分；每题 10 分，共 50 分）}

\SubQ（10 分）试述习近平新时代中国特色社会主义思想创立的时代背景。

\vspace{6pt}
\SubQ（10 分）中国式现代化是如何创造人类文明新形态的？请结合推进中国式现代化必须牢牢把握的重大原则加以论述。

\vspace{6pt}
\SubQ（10 分）为什么说中国共产党的领导是中国特色社会主义最本质的特征和制度的最大优势？

\vspace{6pt}
\SubQ（10 分）如何理解全体人民共同富裕是中国特色社会主义的本质要求？

\vspace{6pt}
\SubQ（10 分）新时代全面深化改革开放为什么是一场深刻革命？请结合全面深化改革总目标加以论述。

\vspace{16pt}
\examrule
\vspace{4pt}
\begin{center}
  {\large\bfseries ——试题到此结束，请认真检查——}
\end{center}
\vspace{4pt}
\examrule

% ============================================================
%  参考答案
% ============================================================
\newpage
\setcounter{bigq}{0}

\begin{center}
  {\LARGE\bfseries 参考答案与评分标准}
\end{center}
\vspace{8pt}
\examrule

% -------- 第一大题参考答案 --------
\BigQ{单项选择题参考答案}

\begin{answerbox}

\textbf{答案速查：}

\begin{center}
\begin{tabular}{|c|c|c|c|c|c|c|c|c|c|}
\hline
题号 & 1 & 2 & 3 & 4 & 5 & 6 & 7 & 8 & 9 \\
\hline
答案 & D & C & B & D & A & C & A & B & B \\
\hline
\end{tabular}

\vspace{4pt}

\begin{tabular}{|c|c|c|c|c|c|c|c|c|c|}
\hline
题号 & 10 & 11 & 12 & 13 & 14 & 15 & 16 & 17 & 18 \\
\hline
答案 & A & C & D & B & B & D & B & B & B \\
\hline
\end{tabular}

\vspace{4pt}

\begin{tabular}{|c|c|c|c|c|c|c|c|c|c|c|c|c|}
\hline
题号 & 19 & 20 & 21 & 22 & 23 & 24 & 25 & 26 & 27 & 28 & 29 & 30 \\
\hline
答案 & B & C & B & A & B & B & B & D & C & A & B & B \\
\hline
\end{tabular}
\end{center}

\vspace{8pt}
\textbf{简要解析：}

\begin{enumerate}[itemsep=2pt, topsep=4pt, leftmargin=1.6em]
\item D项"全面建成小康社会的历史任务已经完成"属于新时代伟大变革的里程碑成就，不属于思想创立的"时代背景"五个方面。
\item "十个明确"是核心关键组成部分、"四梁八柱"；"十四个坚持"是基本方略；"六个必须坚持"是世界观和方法论；"十三个方面成就"展示实践成果，注意区分。
\item "两个结合"中，马克思主义基本原理是"魂脉"，中华优秀传统文化是"根脉"，不能颠倒。
\item 课堂明确强调思想的历史地位只讲"新的飞跃"，没有官方"第三次历史性飞跃"的提法，答题不应自行添加。
\item "两个确立"与"两个维护"是易混点：B项是"两个维护"的内容，不要混为一谈。
\item 科学社会主义基本原则首次在《共产党宣言》中全面阐述，后在《资本论》及其手稿中充分展开，在《哥达纲领批判》《反杜林论》中进一步阐发。
\item 新时代是我国发展新的历史方位，是导论之后第一章的唯一核心考点。
\item 变化的是社会主要矛盾，没有变的是社会主义初级阶段基本国情和发展中国家国际地位，注意"变"与"不变"对应内容不能张冠李戴。
\item 道路问题的实质在于一种主义能否解决国家面临的历史性课题，这是理解"道路决定命运"的核心逻辑。
\item "两个奇迹"指经济快速发展奇迹和社会长期稳定奇迹。
\item 中国式现代化展现的是不同于西方现代化模式的新图景，C项"完全照搬西方模式"与其本质相悖。
\item 五大重大原则为：党的领导、中国道路、人民中心、改革开放、斗争精神，不包含产业政策层面的具体提法。
\item 团结奋斗是中国人民创造历史伟业的必由之路，是"五个必由之路"之一。
\item 中国共产党的领导是中国最大的国情，也是中国特色社会主义最本质的特征。
\item 党的自身优势"四力"中，政治领导力是马克思主义政党第一位的能力。
\item 党中央集中统一领导是党的领导的最高原则，与民主集中制内在一致而不矛盾。
\item 人民立场是中国共产党的根本政治立场，这是第四章第一考点。
\item 共同富裕要坚持有底线、渐进式推进，不能理解为整齐划一、同步实现。
\item 全面深化改革总目标是完善和发展中国特色社会主义制度、推进国家治理体系和治理能力现代化，两句话缺一不可。
\item 方法论中"主线—重点—突破口"层次分明：制度建设为主线，经济体制改革为重点，重要领域和关键环节改革为突破口。
\item 协调发展注重解决发展不平衡问题；创新解决动力问题，绿色解决人与自然和谐问题，共享解决公平正义问题。
\item "两个毫不动摇"是社会主义基本经济制度中所有制结构的基础性表述。
\item 三个"第一"的提法要与教育优先发展、科技自立自强、人才引领驱动一体推进的实践要求对应记忆。
\item 易错点：是"成果民主"而非"结果民主"，两字之差，表述不同，需特别注意。
\item 人民代表大会制度是我国的根本政治制度；其他三项属于基本政治制度，注意层级区分。
\item 文化自信是更基础、更广泛、更深沉的力量，是四个自信中的独特地位。
\item "育新人"对应立德树人、以文化人，培养能够担当民族复兴大任的时代新人。
\item 初次分配发挥基础性作用，再分配发挥调节作用，需健全第三次分配机制，三者层次不同。
\item 党的十八大把生态文明建设纳入中国特色社会主义"五位一体"总体布局。
\item "绿水青山就是金山银山"揭示了保护生态环境就是保护生产力、改善生态环境就是发展生产力的道理，体现发展和保护协同共生。
\end{enumerate}

\end{answerbox}

% -------- 第二大题参考答案 --------
\BigQ{多项选择题参考答案}

\begin{answerbox}

\textbf{答案速查：}

\begin{center}
\begin{tabular}{|c|c|c|c|c|c|}
\hline
题号 & 1 & 2 & 3 & 4 & 5 \\
\hline
答案 & ABCDE & BC & ABCDE & ABCDE & ABCD \\
\hline
\end{tabular}

\vspace{4pt}

\begin{tabular}{|c|c|c|c|c|c|}
\hline
题号 & 6 & 7 & 8 & 9 & 10 \\
\hline
答案 & ABCDE & ABCD & ABC & ABC & ABCDE \\
\hline
\end{tabular}
\end{center}

\vspace{8pt}
\textbf{简要解析：}

\begin{enumerate}[itemsep=2pt, topsep=4pt, leftmargin=1.6em]
\item 时代背景共五个方面，五项均为官方表述，应全选。
\item D、E为干扰项：社会主要矛盾的变化没有改变我们对社会主义所处历史阶段的判断，我国也尚未全面实现社会主义现代化。
\item 中国式现代化的中国特色共五个方面，五项均正确，应全选。
\item 推进中国式现代化必须牢牢把握的重大原则共五个方面，应结合"四项基本原则"理解记忆，五项均正确。
\item E"经济调控力"为干扰项，党的自身优势"四力"仅包括政治领导力、思想引领力、群众组织力、社会号召力。
\item 新发展理念五大内容为创新、协调、绿色、开放、共享，五项均正确。
\item E"结果民主和形式民主的统一"为干扰项，混淆了"成果民主"概念，全过程人民民主的"四个相统一"为A、B、C、D。
\item D"人民代表大会制度"是根本政治制度而非基本政治制度，E"民主集中制"是党的组织原则，二者均为干扰项。
\item 中国特色社会主义文化由中华优秀传统文化、革命文化、社会主义先进文化三部分构成，D、E为干扰项。
\item 原文共列出六对重大关系（含"自立自强与对外开放的关系"），本题所列五项均属于其中正确内容，应全选。
\end{enumerate}

\end{answerbox}

% -------- 第三大题参考答案 --------
\BigQ{简答题参考答案}

\begin{answerbox}

\textbf{第（1）题　试述习近平新时代中国特色社会主义思想创立的时代背景。（10 分）}

习近平新时代中国特色社会主义思想是在科学把握和积极应对一系列重大时代课题中孕育、形成和发展起来的，其创立的时代背景主要包括以下五个方面：

第一，世界百年未有之大变局加速演进。国际力量对比深刻调整，世界格局和国际体系正经历深刻变化，这为这一思想的形成提供了宏阔的国际背景。

第二，中华民族伟大复兴进入关键时期。党和人民事业站在新的历史起点上，实现中国梦既面临难得的历史机遇，也面对一系列重大风险挑战，迫切需要科学理论指引方向。

第三，中国式现代化全面推进拓展。在推进中国式现代化的进程中不断遇到深层次矛盾和问题，需要理论创新破解新的发展难题、开创现代化新路。

第四，科学社会主义在21世纪的中国焕发新的蓬勃生机。中国特色社会主义事业取得的历史性成就，使科学社会主义在21世纪的中国展现出强大生命力。

第五，中国共产党在革命性锻造中更加坚强有力。党坚持自我革命，党的执政能力和领导水平不断提升，为应对各种风险挑战提供了坚强政治保证。

正是在以上五个方面历史条件的综合作用下，习近平新时代中国特色社会主义思想应运而生，成为当代中国马克思主义、二十一世纪马克思主义。

\textit{评分标准：总论点1分；五个方面每点1.6分（共8分），仅列要点未展开论述每点扣0.5分；总结1分。}

\vspace{6pt}
\textbf{第（2）题　中国式现代化是如何创造人类文明新形态的？请结合推进中国式现代化必须牢牢把握的重大原则加以论述。（10 分）}

中国式现代化创造人类文明新形态，主要体现在以下三个方面：

第一，提供了全新的现代化路径，超越了西方现代化模式，展现出不同于西方现代化的新图景，打破了西方现代化道路的单一模式。

第二，为广大发展中国家走向现代化提供了新的选择，打破了"现代化就是西方化"的迷思，证明现代化道路是多样的。

第三，推进中国式现代化必须牢牢把握重大原则，这是创造人类文明新形态的根本保证。这些原则包括：坚持和加强党的全面领导，这是根本政治保证；坚持中国特色社会主义道路，这是必由之路；坚持以人民为中心的发展思想，这是根本价值取向；坚持深化改革开放，这是动力源泉；坚持发扬斗争精神，这是精神支撑。

总之，中国式现代化深深植根于中华优秀传统文化，体现科学社会主义的先进本质，借鉴吸收一切人类优秀文明成果，代表人类文明进步的发展方向，是人类文明新形态的生动体现。

\textit{评分标准：三方面表现每点2分（共6分）；五大原则结合展开2分；总结2分。}

\vspace{6pt}
\textbf{第（3）题　为什么说中国共产党的领导是中国特色社会主义最本质的特征和制度的最大优势？（10 分）}

中国共产党的领导是中国特色社会主义最本质的特征和制度的最大优势，原因主要在于：

第一，中国最大的国情就是中国共产党的领导。中国共产党是最高政治领导力量，党的领导决定着中国特色社会主义的根本性质和发展方向。

第二，党的领导制度是我国的根本领导制度，是中国特色社会主义制度体系的核心，是国家治理体系和治理能力现代化的关键，发挥着提纲挈领、无可替代的作用。

第三，中国共产党的自身优势是中国特色社会主义制度优势的主要来源，体现为政治领导力（马克思主义政党第一位的能力）、思想引领力、群众组织力、社会号召力这"四力"。

第四，从历史和现实看，党的领导是做好党和国家各项工作的根本保证，是战胜一切困难和风险的"定海神针"；新时代党和国家事业取得历史性成就、发生历史性变革，最根本的原因在于党中央的坚强领导和科学理论指引。

第五，坚持党的领导，必须维护党中央权威和集中统一领导，健全党的全面领导制度，但党的领导是全面的、系统的、整体的，并不意味着事事都要由党具体包办。

\textit{评分标准：最大国情判断2分；制度优势来源2分；"四力"展开2分；历史现实依据2分；正确理解党的领导方式2分。}

\vspace{6pt}
\textbf{第（4）题　如何理解全体人民共同富裕是中国特色社会主义的本质要求？（10 分）}

全体人民共同富裕是中国特色社会主义的本质要求，理解这一论断应把握以下几个层面：

第一，性质定位上，共同富裕是中国特色社会主义的本质要求，也是中国式现代化的重要特征，这从根本上区别于以资本为中心的西方现代化。

第二，政治意义上，实现共同富裕不仅是经济问题，而且是关系党的执政基础的重大政治问题，直接关乎人心向背和长期执政基础。

第三，推进方式上，要从全局角度把握共同富裕，正确处理先富和共富的关系，允许一部分人先富起来，同时积极推动先富带后富；必须坚持有底线、渐进式推进，而不能理解为整齐划一、同步实现。

第四，原则与思路上，要把握好鼓励勤劳创新致富、坚持基本经济制度、尽力而为量力而行、坚持循序渐进的原则；构建初次分配、再分配、第三次分配协调配套的制度安排，形成中间大、两头小的橄榄型分配结构。

第五，推动全体人民共同富裕与促进人的全面发展是高度统一的，二者相互促进、相辅相成。

\textit{评分标准：性质定位2分；政治意义2分；推进方式2分（须点出"有底线、渐进式"）；原则思路2分；总结2分。}

\vspace{6pt}
\textbf{第（5）题　新时代全面深化改革开放为什么是一场深刻革命？请结合全面深化改革总目标加以论述。（10 分）}

新时代全面深化改革开放之所以是一场深刻革命，主要体现在：

第一，就其艰巨性、复杂性和系统性来说，全面深化改革开放是涉险滩、闯难关的变革，触及深层次体制机制和利益格局调整。

第二，它是一场思想理论的深刻变革、改革组织方式的深刻变革、国家制度和治理体系的深刻变革、人民广泛参与的深刻变革，涉及范围之广、程度之深前所未有。

第三，它是在多年改革开放基础上的深化，是一场全面、系统、整体的制度创新，涵盖经济、政治、文化、社会、生态文明、党的建设等各领域改革，使中国特色社会主义制度的优越性得到更好发挥。

第四，必须紧紧把握全面深化改革总目标，即完善和发展中国特色社会主义制度、推进国家治理体系和治理能力现代化。这一总目标是有机整体："完善和发展中国特色社会主义制度"规定了改革的根本方向，要求无论怎么改都必须坚持中国共产党领导、坚持中国特色社会主义；"推进国家治理体系和治理能力现代化"明确了改革的鲜明指向，要求通过改革把制度优势转化为国家治理效能。

第五，新时代的改革开放是全方位、深层次、根本性的，取得的成就是历史性、革命性、开创性的，这正是其"深刻革命"性质的集中体现。

\textit{评分标准：深刻革命的艰巨性复杂性判断2分；四个深刻变革2分；全面系统整体的制度创新2分；总目标两句话完整阐述2分；成就性质总结2分。}

\end{answerbox}

\vspace{12pt}
\examrule
\begin{center}
  \textit{参考答案仅供参考，评分时应酌情给分，论述题言之有理、要点齐全且能展开论述可得相应分数。}
\end{center}

\end{document}
